\section{Discussion}\label{sec:discussion}

\textbf{What this shows.} Across three tasks with different objectives, loss functions, and
evaluation metrics, we could not detect a reliable benefit from this system's
similarity-based audio graph beyond what its own node features already provide -- though the
evidence differs by task, not uniform. Task 2's pooling-matched comparison favors the
edge-free encoder directly. Task 3's own standalone comparison is mixed and pooling-dependent
and is \emph{not} the evidence we rely on there; its branch-zeroing and complementarity
diagnostics, both on held-out test data, establish two different, narrower things: branch-
zeroing shows the trained fusion model relies only slightly on the graph branch (topology and
node features together, since zeroing removes both -- not topology-isolating), and
complementarity shows the independently-trained graph-only model detects almost nothing text
misses (a claim about that model, not the fusion model). Together, in this evidence base,
they point to little demonstrated graph value for Task 3 -- a follow-up check later in this
section found the opposite in a second implementation, unresolved which reading generalizes.
Task 4 shows no reliable difference across three seeds. Tasks 2 and 4 rely on
width/depth-matched, lower-parameter edge-free controls; Task 3's such control is
validation-only, so its test-set evidence comes from branch-zeroing and complementarity
instead, neither an edge-free-encoder comparison. The three findings are not built from one
identical comparison type, stated at each step rather than only here.

\textbf{What this does not show.} It does not show that graph neural networks are a poor
choice for music understanding in general, or that no graph construction could help this
kind of system -- only that \emph{this specific} feature-induced similarity graph, applied to
\emph{these} node features, on \emph{these} datasets, gave the model little it did not
already have. A graph built from a genuinely external relation -- not recomputable from the
same node features, such as annotated chord progressions, beat structure, or section
repetition -- is not ruled out and remains untested here.

\textbf{Why this might be happening, and a genuine complication.} The clearest candidate
explanation is structural: a similarity threshold applied to the same features used as node
content produces an adjacency matrix that is, information-theoretically, redundant with those
features (Section~\ref{sec:cross-task}). A secondary, speculative possibility: the limited
topology present is not fully exploited by the trained GraphSAGE model -- the broader GNN
literature separately finds trained message-passing layers can under-use structure through
over-smoothing \cite{li2018oversmoothing}, and a fixed propagation step can match or exceed a
trained GCN-style layer \cite{wu2019sgc}. Table \ref{tab:task2-full}'s exact 2-hop result
(0.6630, beating both edge-free numbers and GraphSAGE, not the CNN baseline at 0.7654) is a
genuine complication for a "topology never helps" reading, reported rather than explained
away -- but not a clean ablation: the 2-hop pipeline uses full-batch Adam, 200 epochs, no
dropout, while GraphSAGE's script does not pass its configured dropout through and trains at
the default 0.2 (Section~\ref{sec:protocol} discloses this bug precisely), with mini-batches
and early stopping (both share weight decay $10^{-4}$). These recipe differences, not only
whether aggregation is learned, could account for some or all of the gap, so we treat it as
exploratory -- a cleaner, topology-only result exists in a companion audit
\cite{tismir_audit} (Section~\ref{sec:results-t2}), a same-architecture
real/shuffled/random-edge factorial, but we flag rather than lean on it: it is an
unpublished, anonymized-for-review manuscript with no stable public identifier at writing,
cited rather than reproduced, its methods and tables not included alongside this paper or
independently inspectable here. We lack an equivalent 2-hop control for Tasks 3/4, and this
result is single-run despite sharing GraphSAGE's seed (42); porting its script into this
repository (\texttt{task2\_exact\_2hop\_control.py}) also found it not exactly reproducible
across environments -- an independent rerun with identical code, seed, and data reached
0.6551, not 0.6630, plausibly floating-point non-associativity in the accumulation step. Both
values still clearly exceed GraphSAGE (0.5976) and both edge-free readings (0.6049, 0.6425),
not changing the qualitative finding, but a further reason to treat this result as
exploratory. We do not treat it as overturning the paper's pattern -- but it keeps this paper
from concluding the deployed encoders extract everything the node features offer.

\textbf{Limitations.} The caption-sanitization control (69.1\% lexical exposure, 84.3\%
retention after retraining, Section~\ref{sec:data}) was run in an internal research directory
outside the two implementations this paper draws from; the stored numbers were directly
re-verified against their result files, but the training script itself is not included
alongside this paper, a reproducibility gap disclosed here rather than left implicit. Task 2
and Task 3's main comparisons use a single training seed each; only Task 4's headline
comparison repeats across seeds. GTZAN carries a small, previously undocumented amount of
exact-duplicate, cross-split leakage (Section~\ref{sec:data}); MusicCaps was checked clean at
the ID level but not audio-content level, so duplicate underlying recordings, if any, would
not have been caught. Neither dataset's artist-level leakage could be checked, since neither
ships usable artist metadata. This paper uses the original, denser graph construction
throughout (Section~\ref{sec:data}); a companion audit of the same system repairs this
construction and finds the same qualitative pattern on repaired graphs \cite{tismir_audit},
but this paper's own four-task comparisons were not rerun on it. Retrieval evaluation here
treats only the one official caption as correct; a semantically reasonable but unpaired clip
would be scored as an error. We checked this directly on 20 retrieval failures
(Section~\ref{sec:results-t4}) and did not find evidence of a large hidden population of such
cases -- most failures were genuine mismatches, the true clip ranked far outside the top 10.
This used only one checkpoint, 20 non-randomly-selected examples, and one unblinded rater
with no predefined rubric, so it does not rule out the effect entirely, only that it is not
large enough to be obvious on a direct look.

\textbf{A follow-up check that complicates the Task 3 finding, and what we do and do not know
about why.} We reran Task 3's branch-zeroing and complementarity diagnostics in a second,
independent implementation of this same architecture family -- the one Task 4's own numbers
are drawn from -- and got a different answer. Removing the graph branch there costs 0.12
macro-F1 at that implementation's own epoch budget and 0.1672 when retrained at the shorter
budget used before (a single archived run; we did not force deterministic training, so, as
with Table~\ref{tab:task2-full}'s 2-hop result, an exact rerun could land differently; no
additional runs were archived to quantify that); complementarity finds 169 of 359 rescued
cells (47.1\%) genuine at the longer budget and 142 of 427 (33.3\%) at the shorter, both well
above the 3 of 508 (0.6\%) reported above. Both readings are real measurements, needing an
explanation rather than a choice of which to believe. The larger effect persists at the
shorter budget, so the longer budget alone does not explain the disagreement. One plausible
contributor is class weighting: unlike this section's complementarity checkpoint (trained
without positive-class weighting, collapsed to predicting a tag present for only 5 of 50),
the second implementation weights all four fusion modes, including \textit{gnn\_only},
throughout. We have not isolated weighting as the cause: the implementations also differ in
split size, tag vocabulary (50 here, 30 there), and segment length, and no
weighted-vs-unweighted comparison was run within one implementation to separate weighting
from these. In the second implementation, branch-zeroing found substantially greater
fusion-model reliance and complementarity substantially more genuine rescue by its own
graph-only model; neither isolates topology specifically, since branch-zeroing does not
separate topology from node information and complementarity is a statement about a separate
model, not the fusion model's behavior. We report this as an open, only partly explained
disagreement, not a correction to the numbers above and not evidence the graph helps once
weighted correctly.

\textbf{Why richer relations might behave differently.} A relation that is not a function of
the node features -- for example, which segments belong to the same repeated section, or
which two chords typically follow one another -- would fall outside the redundancy argument
in Section~\ref{sec:cross-task} by construction, since it encodes something the feature
vectors alone cannot recover. Testing this is a natural next step, not attempted here beyond
the informal argument above.
